\section{ماژول cachenotes\_compressor (CacheNotes Compressor)}
\label{sec:module_cachenotes_compressor}

\subsection{معرفی و انگیزه}
این ماژول در چارچوب پژوهش حاضر پیاده‌سازی و ارزیابی شده و فاقد مقاله مستقل جداگانه در ادبیات تحقیق است.

در پردازش دنباله‌های طولانی توسط ترنسفورمرها، رشد خطی حجم حافظه پنهان کلید و مقدار (\lr{KV Cache}) به همراه نیاز به جابجایی وزن‌های حجیم شبکه، تنگنای اصلی پهنای باند حافظه (\lr{Memory Bandwidth Bottleneck}) را پدید می‌آورد. اکثر روش‌های پیشین فشرده‌سازی و هرس، کاملاً مستقل از ماهیت وظیفه پیش‌رو عمل کرده و فرض می‌کنند اهمیت کانال‌ها برای تمامی وظایف زبانی یکسان است؛ حال آنکه نیازمندی‌های استخراج ویژگی در وظایفی چون خلاصه‌سازی (\lr{Summarization})، پاسخ به پرسش (\lr{QA})، یا کدنویسی به شدت با یکدیگر متمایز هستند.

چارچوب فشرده‌سازی \lr{CacheNotes} برای ایجاد فشرده‌سازی آگاه از وظیفه (\lr{Task-Aware Compression}) معرفی شده است. این چارچوب با تولید توکن‌های پیوسته برنامه‌ریزی فشرده‌سازی (\lr{Compression-Planning-Tokens - CPTs})، اطلاعات معنایی وظیفه را به بردار اهمیت پیوسته کانال‌ها نگاشت می‌کند. سپس این بردار هم در فشرده‌سازی استاتیک وزن‌ها بر پایه تجزیه مقادیر منفرد هدایت‌شده (\lr{CPT-Guided SVD}) و هم در هرس پویای حافظه پنهان \lr{KV Cache} در زمان اجرا به کار گرفته می‌شود.

\subsection{مبانی نظری و فرمول‌بندی ریاضی}
فرآیند فشرده‌سازی در \lr{CacheNotes} از دو مولفه ریاضی وابسته تشکیل شده است:
\begin{enumerate}
    \item \textbf{تولیدکننده پیوسته توکن برنامه‌ریزی فشرده‌سازی (\lr{CPT Generator}):} با دریافت توصیف متنی وظیفه $T_{\text{task}}$، ابتدا امبدینگ توکن‌ها توسط یک شبکه رمزگذار به بردار حالت $\mathbf{h} \in \mathbb{R}^{d_h}$ تبدیل می‌شود. برای تولید اهمیت در هر بعد دلخواه $d$ (حتی ابعاد پیوسته)، موقعیت کانال‌ها به صورت فرکانسی در قالب مختصات \lr{NeRF} کدگذاری می‌شود:
    \begin{equation}
    \gamma(c) = \left[ \sin(2^0 \pi c), \cos(2^0 \pi c), \dots, \sin(2^{F-1} \pi c), \cos(2^{F-1} \pi c) \right], \quad c = \frac{j}{d}, \quad j \in \{1, \dots, d\}
    \end{equation}
    بردار اهمیت نهایی کانال‌ها توسط شبکه پروجکشن بر روی بازه پیوسته $[0.1, 2.0]$ نگاشته می‌شود:
    \begin{equation}
    \mathbf{s} = \operatorname{Sigmoid}(\text{MLP}([\mathbf{h}; \gamma(\mathbf{c})])) \cdot 1.9 + 0.1 \in \mathbb{R}^d
    \end{equation}
    \item \textbf{تجزیه کم‌مرتبه هدایت‌شده (\lr{CPT-Guided SVD}):} ماتریس وزن هدف $\mathbf{W} \in \mathbb{R}^{d_{\text{out}} \times d_{\text{in}}}$ با ضرب در قطر بردار اهمیت ورودی یا خروجی پیش‌شرط‌سازی می‌گردد:
    \begin{equation}
    \tilde{\mathbf{W}} = \mathbf{W} \operatorname{diag}(\mathbf{s})
    \end{equation}
    سپس بر روی $\tilde{\mathbf{W}}$ تجزیه مقادیر منفرد برشی اعمال شده و فاکتورهای کم‌مرتبه پس از بازگردانی مقیاس تشکیل می‌گردند:
    \begin{equation}
    \tilde{\mathbf{W}} \approx \mathbf{U}_r \mathbf{\Sigma}_r \mathbf{V}_r^\top \implies \hat{\mathbf{W}} = \mathbf{U}_r \mathbf{\Sigma}_r \mathbf{V}_r^\top \operatorname{diag}(\mathbf{s})^{-1}
    \end{equation}
    \item \textbf{فشرده‌سازی پویای حافظه پنهان (\lr{KV Cache Compression}):} در زمان استنتاج، توکن‌های حافظه پنهان کلید و مقدار با ضرب در وزن‌های اهمیت $\mathbf{s}$ ارزیابی شده و کم‌اهمیت‌ترین توکن‌ها تا رسیدن به بودجه فشرده‌سازی هرس می‌شوند.
\end{enumerate}

\subsection{معماری و پیاده‌سازی}
کد منبع این ماژول در فایل \lr{neuralops/compression\_engine/cachenotes\_compressor.py} مستقر بوده و شامل بخش‌های زیر است:
\begin{itemize}
    \item \textbf{کلاس \lr{CPTGenerator}:} شبکه عصبی تفکیک‌پذیر مجهز به \lr{EmbeddingBag} و لایه‌های پروجکشن فرکانسی که امکان ارزیابی وظیفه در هر بعد دلخواه $d$ (مانند ۶۴، ۱۲۸، ۷۶۸، ۱۰۲۴ یا ۴۰۹۶) را بدون محدودیت معماری فراهم می‌سازد.
    \item \textbf{کلاس \lr{CompressionPlanningTokens}:} سیستم پشتیبان مبتنی بر کلیدواژه‌های ایستا که در صورت عدم آموزش شبکه CPT، به طور خودکار به عنوان جایگزین پایدار فعال می‌شود.
    \item \textbf{کلاس \lr{CacheNotesLoRALinear}:} لایه جایگزین فاکتورگیری‌شده که وزن فشرده را ثابت نگه داشته و شاخه کم‌مرتبه \lr{LoRA} را جهت بازیابی دقیق تعبیه می‌کند.
    \item \textbf{کلاس \lr{CacheNotesCompressor}:} کلاس اصلی مدیریت فرآیند فشرده‌سازی، مدیریت وظایف کالیبراسیون و کالیبراسیون مجدد لایه‌نرمال.
\end{itemize}

\subsection{نتایج تجربی}
نتایج ارزیابی آزمایشگاهی فشرده‌ساز \lr{cachenotes\_compressor} بر روی مدل‌های مرجع در جدول \ref{tab:benchmark_cachenotes_compressor} خلاصه شده است:

\begin{table}[H]
\centering
\small
\caption{نتایج بنچ‌مارک سخت‌افزاری و زبانی فشرده‌ساز \lr{cachenotes\_compressor} روی مدل‌های پایه}
\label{tab:benchmark_cachenotes_compressor}
\begin{tabular}{lccccc}
\toprule
\textbf{مدل پایه \lr{(Base Model)}} & \textbf{دقت \lr{(Accuracy)}} & \textbf{شاخص \lr{F1}} & \textbf{حافظه \lr{(Peak VRAM)}} & \textbf{توان \lr{(Throughput)}} & \textbf{وضعیت} \\
\midrule
\lr{DistilBERT (SST-2)} & \lr{90.94\%} & \lr{0.9127} & \lr{368.29 MB} & \lr{133,508.96 tok/s} & \lr{PASS} \\
\lr{DistilBERT (MNLI)} & \lr{81.70\%} & \lr{0.8162} & \lr{456.73 MB} & \lr{193,047.72 tok/s} & \lr{PASS} \\
\lr{RoBERTa (SST-2)} & \lr{92.55\%} & \lr{0.9274} & \lr{514.29 MB} & \lr{12,830.05 tok/s} & \lr{PASS} \\
\lr{GPT-2 (AGNews)} & \lr{81.30\%} & \lr{0.8080} & \lr{743.20 MB} & \lr{40,222.56 tok/s} & \lr{PASS} \\
\lr{Qwen} & -- & -- & -- & -- & \lr{NOT\_RUN} \\
\bottomrule
\end{tabular}
\end{table}

بر پایه نتایج تجربی، فشرده‌ساز \lr{cachenotes\_compressor} عملکرد بسیار متعادل و قدرتمندی را در تمامی مدل‌ها به نمایش گذاشته است. در مدل \lr{RoBERTa} به دقت عالی \lr{92.55\%} دست یافته و در مدل \lr{DistilBERT} دقت \lr{90.94\%} را روی SST-2 با توان بالای \lr{133,508} توکن بر ثانیه ثبت نموده است. علاوه بر این، بر روی مدل دیکودری \lr{GPT-2} بر روی دادگان AGNews دقت بالای \lr{81.30\%} با توان \lr{40,222} توکن بر ثانیه به دست آمده و وضعیت \lr{PASS} را احراز کرده است. آزمون روی مدل \lr{Qwen} انجام نشده و وضعیت \lr{NOT\_RUN} ثبت شده است.

\subsection{محدودیت‌ها و تحلیل شکست}
تحلیل ساختاری و تجربی این ماژول نکات زیر را مشخص می‌سازد:
\begin{itemize}
    \item \textbf{وابستگی به توصیف کیفی وظیفه:} عملکرد ماژول به توصیف متنی ارائه‌شده در ورودی وابسته است؛ چنانچه توصیف وظیفه مبهم یا کوتاه باشد، بردار اهمیت تولیدشده ممکن است تنوع کافی را ایجاد نکند و ماژول ناگزیر به استفاده از مقادیر پیش‌فرض هیوریستیک گردد.
    \item \textbf{سربار فاز اولیه آموزش مولد \lr{CPT}:} برای دستیابی به حداکثر تطبیق‌پذیری، شبکه \lr{CPTGenerator} نیازمند آموزش پیشینی بر روی دادگان وظایف است؛ اگرچه مکانیزم پشتیبان ایستا مانع از شکست عملکردی می‌شود، اما بهره‌برداری کامل از پتانسیل مدل مشروط به مرحله اولیه کالیبراسیون وظیفه‌محور است.
\end{itemize}
